
%\hl{discuss the difference between material reductions driven by structural load decrease versus structural efficiency increase, and how to tell the difference and which is happening here.}


%\hl{describe which applications are best suited to which design. Maybe do min-LCOE, more expensive but could be min LCOE under some different economic scenario, cheaper and 10kW PBE, and cheaper and 1kW PBE.}

%This design is recommended for applications that are highly sensitive to the cost of energy and where energy storage is relatively cheap and abundant or there are dispatchable loads. This encompasses utility-scale installations and integrated generation-storage concepts such as a WEC coupled with offshore hydrogen production or pumped hydro. 

%On the opposite end of the spectrum, a second design is the least variable design, with an LCOE of \$1.00/kWh and merely 58\% power variation. This design is recommended for applications where the cost of energy is less critical, and storage is extremely limited. Powering the Blue Economy (PBE) applications such as aquaculture farms and AUV charging likely meet these criteria. 

%Finally, we recommend the balanced design for intermediate applications, where energy storage is available but not abundant, and cost is important but not the bottom line. An example is a grid-connected application for island communities. This balanced design has an LCOE of \$0.18/kWh and a power coefficient of variation equal to 100\%. 
%Further insight on the power profiles of these recommended designs compared to the baseline can be uncovered through probability distribution plots, shown in \Cref{fig:power-dist}. The minimum variation design is seen to be tuned to produce most of its power from moderate sea states only ($\sim$60-300 kW), whereas the minimum LCOE design has a much broader spectrum, collecting significant power both below 10~kW and above 1~MW. The balanced design falls in the middle. %As expected, the minimum LCOE design operates at high powers the most, spending 50\% of the time above 500~kW, compared to 40\% and 10\% for the balanced and minimum variation designs respectively. Also as expected, the minimum variation design spends a large amount of time ($\approx$40\%) at low powers below 50~kW. Interestingly, the minimum-LCOE design actually spends more time at low powers than the balanced design. In future work, the power distributions could be optimized more explicitly, incorporating information other than the coefficient of variation as an objective based on grid-level analysis.


%Highly sensitive parameters include \hl{XXX}.
%Insensitive parameters include \hl{XXX.
%Describe which are sensitive to J vs x vs constraint activity.
%Maybe also show dJ/dx and that's important because you need that DV exact or not without affecting your bottom line.}
%\hl{Is it possible to use these sensitivities to demonstrate coupling, and predict the effect of holding some DVs constant?) - during re scrutineering} 

%For these highly sensitive parameters, it is worthwhile to perform a global sensitivity analysis to observe any nonlinearity. \Cref{remdo:fig:param-sens-global,remdo:fig:x-opt-sens-global} demonstrate the new optimal objective and optimal design variable values respectively under the re-optimization-based OAT global parameter sensitivity process described in \Cref{sec:single-obj-process}.
% \begin{figure}\label{remdo:fig:param-sens-global}
% \includegraphics[width=\linewidth]{example-image-a}
% \caption{Global OAT parameter sensitivity analysis to optimal objective $J^*$}
% \end{figure}
% \begin{figure}\label{remdo:fig:x-opt-sens-global}
% \includegraphics{example-image-a}
% \caption{Global OAT parameter sensitivity analysis to optimal design $x^*$}
% \end{figure}

%The global sensitivity analysis indicates that \hl{list parameters} have strongly nonlinear sensitivities over the range of interest. During re-scrutineering. \hl{Sentence interpreting any params that aren't location based}. 

%%%%%%%%%%%%%%%%%%%%


%\subsubsection{Price Global Sensitivity}
%We also perform a global sensitivity analysis for price parameters $p_\tau$, $p_P$, $p_M$, and $p_0$ from equation~\eqref{remdo:eq:LCOE-scale}. For $p_0$, if the design-independent (unmodeled) cost goes to zero, \hl{(and maybe change 80\% efficiency into 100\%)}, that's essentially the best LCOE that a WEC of these technical assumptions (steel point absorber with given operational conditions) can get. We can think of the rest as a cost gap \hl{(explain)}. Meanwhile, the global sensitivity analysis for $p_\tau$ gives preliminary insight into the multi-objective problem, recovering a portion of the pareto front between the four variables that influence $LCOE$: $\tau$, $P_{pk,\text{elec}}$, $V_{\text{struct}}$, and $P_{\text{avg},\text{elec}}$. This shows \hl{XXX conclusion}, to be clarified further in the multi-objective optimization.

%%%%%%%%%%%%%%%%%%%%
% \subsubsection{Global Sensitivity to Modeling Assumptions}
% \label{remdo:sec:sensitivity-model-assumptions}
% Beyond continuous parameters, the simulation depends on discrete options governing model behavior: the control scheme (damping or reactive), the presence of generator force saturation, the presence of drag, and the number of dynamic degrees of freedom (float-heave only, or float and spar heave).

%\paragraph{Control Scheme}
%\hl{This analysis is left for during scrutineering.}

%%%
% \paragraph{Force Saturation}
\ifdefined\DISSERTATION
    For the minimum-LCOE design of \Cref{remdo:sec:results-single}, the lowest force-saturation power factor across all sea states in \Cref{remdo:fig:power-matrix-opt} is 
    $P_{sat}/P_{\text{unsat}}=\resultsRE[lowestFmaxFactorMinLCOE]$, at a force limit of $F_{\text{max}}/F_{\text{unsat}}=\resultsRE[lowestFsatMinLCOE]$ in that sea state.
    Across all sea states, the optimal force-saturated design sacrifices only 
    \resultsRE[powerLossForceSatMinLCOE]~of average power (about \resultsRE[pctPowerLossForceSatMinLCOE]) while cutting PTO cost by \resultsRE[pctPTOSavingsForceSatMinLCOE],
    for an LCOE reduction of \resultsRE[pctImproveLCOEForceSatMinLCOE]~versus the same design without saturation, which in any case violates structural constraints due to the higher force.
    Re-optimizing to satisfy constraints without allowing force saturation yields a
    \resultsRE[pctPowerGainForceSatOptMinLCOE]~power gain,
    \resultsRE[pctPTOCostIncreaseForceSatOptMinLCOE]~higher PTO cost,
    \resultsRE[pctStructCostIncreaseForceSatOptMinLCOE]~higher structural cost, and
    \resultsRE[pctLCOEIncreaseForceSatOptMinLCOE]~higher LCOE than the optimization with saturation.
\else
    % For the minimum-LCOE design, the optimal force-saturated solution sacrifices only \resultsRE[powerLossForceSatMinLCOE]~of average power (about \resultsRE[pctPowerLossForceSatMinLCOE]) 
    % while cutting PTO cost by \resultsRE[pctPTOSavingsForceSatMinLCOE], for an LCOE reduction of \resultsRE[pctImproveLCOEForceSatMinLCOE]~versus the unsaturated equivalent (which violates structural constraints).
    % Re-optimizing to avoid saturation while satisfying constraints gives \resultsRE[pctLCOEIncreaseForceSatOptMinLCOE]~higher LCOE, confirming that force saturation is advantageous.
\fi
%This corresponds almost exactly with the relationship expected in the purely resistive Thevenin impedance case in \citep{mccabe_force-limited_2024}:
%\begin{equation}
%   \frac{ P_{sat}}{P_{\text{unsat}}}=1-\left(\frac{F_{\text{max}}}{F_{\text{unsat}}}-1\right)^2
%\end{equation}
% see notebook p46 5/16/25
% fixme: what about the factor of 4/pi in the force ratio? Then the quadratic equation wouldn't apply, it would be more complicated (alpha non zero case) which casts doubt on my conclusion that alpha=0 emerges as optimal. Would need to compute alpha directly I think? I also think that the assumption of damping control means my Im(Gamma)=0 assumption in the notebook is false. Also Psat/Punsat is the the same as Psat/Pmax. So leaving this sentence out.

%%%
% \paragraph{Drag}
\ifdefined\DISSERTATION
    As \Cref{remdo:sec:sensitivities-local} notes, the low sensitivity to drag coefficient suggests that modeling the drag nonlinearity may not be essential.
    This is surprising, since without drag the resonant amplitude is significantly higher, which should matter given the active slamming constraint.
    One reason for the discrepancy is that the drag sensitivity was local only; a re-optimization at $C_d=0$ would likely give different results.
\else
    % The low local sensitivity to drag coefficient (\Cref{remdo:sec:sensitivities-local}) suggests modeling the drag nonlinearity may be inessential.
    % This is surprising given the active slamming constraint, though a global re-optimization at $C_d=0$ is needed to confirm.
\fi
%\hl{Performing this re-optimization is left for during scrutineering.}

%\paragraph{Degrees of Freedom}
%\hl{This analysis is left for during scrutineering.}


%An additional point of note on the Pareto front of \Cref{remdo:fig:pareto} is the LCOE and coefficient of variation of a typical solar array plotted for comparison (data from \citep{borettiSolar}). Even the cheapest WEC design is still almost an order of magnitude more expensive than solar, indicative of the immaturity of WEC technology and the need for further innovation and cost reduction. On the bright side, many of the designs achieve a lower coefficient of power variation than solar, indicating the potential for wave energy as a relatively steady grid-connected renewable. Finally, it is worth highlighting that the amount of power variation directly affects the required energy storage, so the local cost of storage affects the position along the Pareto front that is optimal for a given application.

%We can plot LCOE contours under different economic assumptions (sweep a couple FCRs, p0s, and opexes) and overlay this to the results already obtained in the global sensitivity analysis of FCR, p0, and opex but now plotted in multiobjective space.\hl{(Explain this better)}



%\Cref{remdo:fig:pareto-epsilon-runtime-sweep} shows the effect of the number of epsilon-constraint seeds on function evaluations and hypervolume.
%\hl{Analysis is left for during scrutineering.}

% \begin{figure}[htbp]
%     \centering
%     \includegraphics[width=\linewidth]{example-image-a}
%     \caption{Effect of number of epsilon-constraint seeds on pareto front function evaluations and hypervolume}
%     \label{remdo:fig:pareto-epsilon-runtime-sweep}
% \end{figure}
